\documentclass[11pt,letterpaper]{article}

\usepackage[margin=0.85in]{geometry}
\usepackage[T1]{fontenc}
\usepackage[utf8]{inputenc}
\usepackage{lmodern}
\usepackage{microtype}
\usepackage{tabularx}
\usepackage{enumitem}
\usepackage{titlesec}
\usepackage[hidelinks]{hyperref}
\usepackage{xurl}
\urlstyle{same}
\usepackage{fancyhdr}
\usepackage{multicol}

\setlength{\parindent}{0pt}
\setlength{\parskip}{0pt}
\setcounter{secnumdepth}{0}

\titleformat{\section}{\Large\scshape\raggedright}{}{0em}{}[\titlerule]
\titlespacing*{\section}{0pt}{1.2ex}{0.8ex}
\titleformat{\subsection}{\bfseries\large}{}{0em}{}
\titlespacing*{\subsection}{0pt}{1.0ex}{0.5ex}

\setlist[itemize]{leftmargin=*, itemsep=0.2em, topsep=0.2em}
\setlist[enumerate]{leftmargin=*, itemsep=0.2em, topsep=0.2em}

\pagestyle{fancy}
\fancyhf{}
\rfoot{\thepage}
\renewcommand{\headrulewidth}{0pt}

\newcommand{\cventry}[3]{%
  \begin{tabularx}{\textwidth}{@{}X r@{}}%
    \textbf{#1} & \textit{#2} \\
  \end{tabularx}\par
  \if\relax\detokenize{#3}\relax\else{\small #3}\par\fi
  \vspace{0.35em}
}

\newcommand{\cvcourse}[3]{%
  \begin{tabularx}{\textwidth}{@{}X >{\raggedleft\arraybackslash}p{0.33\textwidth}@{}}%
    \textbf{#1} & \textit{#2} \\
  \end{tabularx}\par
  {\small #3}\par
  \vspace{0.35em}
}

\begin{document}

\begin{center}
{\huge \textbf{Chris J. Vargo}}\\[2pt]
1B15 Armory Building \textbullet{} University of Colorado Boulder \textbullet{} Boulder, CO 80303\\
\href{mailto:christopher.vargo@colorado.edu}{christopher.vargo@colorado.edu} \textbullet{} (724) 840-2973 \textbullet{} ORCID: 0000-0003-1733-2506\\
{\small Current as of February 1, 2026}
\end{center}
\vspace{-0.3em}

\section{Education}
\cventry{Ph.D., Mass Communication}{May 2014}{University of North Carolina at Chapel Hill}
\cventry{Master of Arts, Advertising \& Public Relations}{May 2011}{University of Alabama}
\cventry{Bachelor of Arts, Advertising \& Public Relations}{May 2008}{Pennsylvania State University}
\section{Current Academic Appointments}
\cventry{Associate Professor of Advertising}{August 2021 -- Present}{College of Media, Communication \& Information \\ University of Colorado at Boulder}
\cventry{Associate Professor of Information Management and Analytics (Courtesy)}{June 2021 -- Present}{Leeds School of Business \\ University of Colorado at Boulder}
\cventry{Graduate Faculty Member and Instructor}{January 2022 -- Present}{Master of Science in Data Science \\ University of Colorado at Boulder}
\section{Past Academic Appointments}
\cventry{CMCI Director of Master of Science in Business Analytics (MSBA)}{May 2018 -- May 2023}{College of Media, Communication \& Information \\ The University of Colorado at Boulder}
\cventry{Associate Chair of Graduate Studies}{Fall 2022}{Advertising, Public Relations, and Media Design \\ University of Colorado at Boulder}
\cventry{Editor--in--Chief}{December 2018 -- January 2022}{The Agenda Setting Journal}
\cventry{Assistant Professor of Advertising}{August 2016 -- July 2021}{College of Media, Communication \& Information \\ University of Colorado at Boulder}
\cventry{Assistant Professor of Advertising and Public Relations}{August 2014 -- July 2016}{College of Communication \& Information Sciences \\ The University of Alabama}
\cventry{Adjunct Faculty and Graduate Research Assistant}{August 2011 -- May 2014}{School of Media and Journalism \\ The University of North Carolina at Chapel Hill}
\cventry{Adjunct Faculty}{June -- August 2011}{College of Communication \& Information Sciences \\ The University of Alabama}
\cventry{Graduate Teaching Assistant}{August 2009 -- May 2011}{College of Communication \& Information Sciences \\ The University of Alabama}
\section{Research}
\subsection*{Patents}
{\small
\begin{itemize}
\item Vargo, C., Hopp, T. (2021). Systems and Methods for Detecting Social Diversity and Inclusion. Patent Application. No. 63/133,741. CU Reference No.: CU5462B--PPA1. Filed Jan 4th, 2022.
\end{itemize}
}
\subsection*{Books}
{\small
\begin{enumerate}
\item Vargo, C. (2024) The Computational Content Analyst: Using Machine Learning to Classify Media Messages. Routledge. \url{https://www.routledge.com/The-Computational-Content-Analyst-Using-Machine-Learning-to-Classify-Media-Messages/Vargo/p/book/9781032846309}
\item Shaw, D., Minooie, M., Aikat, D., \& Vargo, C. (2019). Agendamelding: News, social media, audiences, and civic community. New York, NY: Peter Lang. \url{https://doi.org/10.3726/b15023}
\item Graham, G., Greenhill, A., Shaw, D., \& Vargo, C. (2015). Content is king: News media management in the digital age. New York, NY: Bloomsbury.
\end{enumerate}
}
\subsection*{Refereed Monographs}
{\small
\begin{enumerate}
\item Amazeen, M., Vargo, C., \& Hopp, T. (2018). Reinforcing attitudes in a gatewatching news era: Individual--level antecedents to sharing fact--checks on social media. Communication Monographs, 86(1), 112--132. \url{https://doi.org/10.1080/03637751.2018.1521984}
\end{enumerate}
}
\subsection*{Refereed Articles}
{\small
\begin{enumerate}
\item Vargo, C. J. , \& Hopp, T. (in press). Incivility on popular politics and news subreddits: An analysis of in-groups, community guidelines and relationships with social media engagement [Accepted October 27, 2025]. Social Science Computer Review.
\item Hopp, T., Ferrucci, P., Vargo, C. J., \& Mays, B. (2025). Mainstream news media trust, countermedia attendance, and political learning. Journalism \& Mass Communication Quarterly, Advance online publication. \url{https://doi.org/10.1177/10776990251329044}
\item Vargo, C. J., \& Minooie, M. (2025). Breaking Agenda Setting Boundaries: A Multi-Dimensional Approach to Understanding Salience of Gun Control in the Polarized Public Sphere. Mass Communication \& Society, 28(3), 530--553. \url{https://doi.org/10.1080/15205436.2024.2333972}
\item Vargo, C. J., Hopp, T., \& Agarwal, P. (2024). Balancing Brand Safety and User Engagement in a Two-Sided Market: An Analysis of Content Monetization on Reddit. Journal of Current Issues \& Research in Advertising, 45(2), 242--256. \url{https://doi.org/10.1080/10641734.2023.2301621}
\item Minooie, M., Taylor, B., \& Vargo, C. (2023). Agendamelding and COVID-19: the dance of horizontal and vertical media in a pandemic. Frontiers in Political Science. (5)2023. \url{https://doi.org/10.3389/fpos.2023.1021855}
\item Hopp, T., Ferrucci, P., Vargo, C., \& Liu, L. (2025). Is Online Textual Political Expression Associated with Political Knowledge? Communication Research, 52(1), 32--60. \url{https://doi.org/10.1177/00936502221113808}
\item Schmierbach, M., McCombs, M., Valenzuela, S., Dearing, J. W., Guo, L., Iyengar, S., Kiousis, S., Kosicki, G. M., Meraz, S., Scheufele, D. A., Stoycheff, E., Vargo, C., Weaver, D. H., \& Willnat, L. W. (2022). Reflections on a Legacy: Thoughts from Scholars about Agenda-Setting Past and Future. Mass Communication and Society, 25(4), 500-527. \url{https://doi.org/10.1080/15205436.2022.2067725}
\item Vargo, C. (2022). Public ``agendamelding'' in the United States: Assessing the relative influence of different types of online news on partisan agendas from 2015 to 2020. Journal of Information Technology \& Politics. 19(3), 284-301. \url{http://dx.doi.org/10.1080/19331681.2021.1972893}
\item Blankenship, J. \& Vargo, C. (2022). The Effect of Corporate Media Ownership on Depth of Local Coverage and Issue Agendas: A Computational Case Study of Six Sinclair TV Station Websites. Electronic News, 15(3-4), 139-158. \url{https://doi.org/10.1177/19312431211043483}
\item Vargo, C. \& Amazeen, M. (2021). Agenda-Cutting Versus Agenda-Building: Does Sponsored Content Influence Corporate News Coverage in U.S. Media? International Journal of Communication 15(2021), 1--22. \url{https://ijoc.org/index.php/ijoc/article/view/17824}
\item Amazeen, M. \& Vargo, C. (2021). Sharing Native Advertising on Twitter: Content Analyses Examining Disclosure Practices and Their Inoculating Influence. Journalism Studies 22(7), 916--933. \url{https://doi.org/10.1080/1461670X.2021.1906298}
\item Hopp, T., Ferrucci, P., Fisher, J., Vargo, C. (2020). Exposure to Difference on Facebook, Trust, and Political Knowledge. Mass Communication \& Society. 23(6), 779--809. \url{https://doi.org/10.1080/15205436.2020.1823002}
\item Guo, L., \& Vargo, C. (2020). Predictors of international news flow: Exploring a networked global media system. Journal of Broadcasting \& Electronic Media. 64(3), 418--437. \url{https://doi.org/10.1080/08838151.2020.1796391}
\item Hopp, T., Ferrucci, P., \& Vargo, C. (2020). Why Do People Share Ideologically Extreme, False, and Misleading Content on Social Media? A Self--Report and Trace Data--Based Analysis of Countermedia Content Dissemination on Facebook and Twitter. Human Communication Research, 46(4), 357--384. \url{https://doi.org/10.1093/hcr/hqz022}
\item Vargo, C., \& Hopp, T. (2020). Associations between advertisement negativity and political advertisement engagement: A computational case study of Russian--linked Facebook and Instagram content. Journalism \& Mass Communication Quarterly. 97(3), 743--761. \url{https://doi.org/10.1177/1077699020911884}
\item Ferrucci, P., Hopp, T., \& Vargo, C. (2019). Civic engagement, social capital, and ideological extremity: Exploring online political engagement and political expression on Facebook. New Media \& Society, 22(6), 1095--1115. \url{https://doi.org/10.1177/1461444819873110}
\item Hopp, T., \& Vargo, C. (2019). Social capital as an inhibitor of online political incivility: An analysis of behavioral patterns among politically active Facebook users. International Journal of Communication, 13(2019), 4313--4333. \url{https://ijoc.org/index.php/ijoc/article/view/10523}
\item Vargo, C., Gangadharbatla, H., \& Hopp, T. (2019). eWOM across channels: Comparing the impact of self--enhancement, positivity bias and vengeance on Facebook and Twitter. International Journal Advertising, 38(8), 1153--1172. \url{https://doi.org/10.1080/02650487.2019.1593720}
\item Vargo, C., \& Hopp, T. (2019). Attention to issues and facts: Assessing the role of need for orientation as a predictor of political news sharing on Facebook. The Agenda--Setting Journal, 3(2), 186--207. \url{https://doi.org/10.1075/asj.18004.var}
\item Guo, L., \& Vargo, C. (2020). ``Fake news'' and emerging online media ecosystem: An integrated intermedia agenda--setting analysis during the 2016 U.S. presidential election. Communication Research, 47(2), 178--200. \url{https://doi.org/10.1177/0093650218777177}
\item Hopp, T., Vargo, C., Dixon, L., \& Thain, N. (2018). Correlating self--report and trace data measures of incivility: A proof of concept. Social Science Computer Review, Special Issue ``Integrating Survey Data and Digital Trace Data.'' 38(5), 584--599. \url{https://doi.org/10.1177/0894439318814241}
\item Vargo, C. (2018). Fifty years of agenda--setting research: New directions and challenges for the theory. The Agenda--Setting Journal, 2(2), 105--123. \url{https://doi.org/10.1075/asj.18023.var}
\item Vargo, C., Guo, L., \& Amazeen, A. (2018). The agenda--setting power of fake news: A big data analysis of the online media landscape from 2014 to 2016. New Media \& Society. 20(5) 2028--2049. \url{https://doi.org/10.1177/1461444817712086}
\item Guo, L., \& Vargo, C. (2017). Global intermedia agenda setting: A big data analysis of international news flow. Journal of Communication, 67(4), 499--520. \url{https://doi.org/10.1111/jcom.12311}
\item Hopp, T., \& Vargo, C. (2017). Does negative campaign advertising stimulate uncivil communication on social media? Measuring audience response using big data. Computers in Human Behavior, 68, 368--377. \url{https://doi.org/10.1016/j.chb.2016.11.034}
\item McGregor, S., \& Vargo, C. (2017). Election--related talk and agenda setting effects on Twitter: A big data analysis of salience transfer at different levels of user participation. The Agenda--Setting Journal, 1(1), 44--63. \url{https://doi.org/10.1075/asj.1.1.05mcg}
\item Tomeny, T., Vargo, C., \& El--Toukhy, S. (2017) Geographic and demographic correlates of autism--related anti--vaccine beliefs on Twitter, 2009--15. Social Science \& Medicine, 191, 168--175. \url{https://doi.org/10.1016/j.socscimed.2017.08.041}
\item Vargo, C., \& Guo, L. (2017). Networks, big data, and intermedia agenda--setting: An analysis of traditional, partisan, and emerging online U.S. news. Journalism \& Mass Communication Quarterly, 94(4), 1031--1055. \url{https://doi.org/10.1177/1077699016679976}
\item Guo, L., Vargo, C., Pan, Z., \& Ding, W. (2016). Big social data analytics in journalism and mass communication: Comparing dictionary--based text analysis and unsupervised topic modeling. Journalism \& Mass Communication Quarterly, 93(2), 332--359. \url{https://doi.org/10.1177/1077699016639231}
\item Kim, Y., Gonzenbach, B., Vargo, C., \& Kim, Y. (2016). First and second levels of intermedia agenda setting: Political advertising, newspapers, and Twitter during the 2012 U.S. presidential election. International Journal of Communication, 10(2016), 4550--4569.
\item Shaw, D., Mousa, I., Vargo, C., Minooie, M., \& Cole, R. (2016). The agenda setting in the digital age: How we use media to monitor civic life and reframe community. Jordan Journal of Social Sciences, 9(1), 125--139.
\item Vargo, C. (2016). Toward a tweet typology: A study of brand message content types and corresponding consumer engagement on social media. Journal of Interactive Advertising, 16(2), 157--168. \url{https://doi.org/10.1080/15252019.2016.1208125}
\item Vargo, C., \& Hopp, T. (2016). Social capital, political polarity, and socioeconomic status as predictors of political incivility on Twitter: A congressional district--level analysis. Social Science Computer Review, 35(1), 10--32. \url{https://doi.org/10.1177/0894439315602858}
\item Guo, L., \& Vargo, C. (2015). The power of message networks: A big--data analysis of the network agenda setting model and issue ownership. Mass Communication \& Society, 18(5), 557--576. \url{https://doi.org/10.1080/15205436.2015.1045300}
\item Vargo, C., Basilaia, E., \& Shaw, D. (2015). Event vs. issue: Twitter reflections of major news, a case study. Studies in Media and Communication, 9, 215--239. \url{https://doi.org/10.1108/S2050-206020150000009009}
\item Shaw, D., \& Vargo, C. (2014). Media and social stability: How audiences create personal community. Journalism Bimonthly [in Chinese], 6(12), 16--24.
\item Vargo, C., Guo, L., McCombs, M., \& Shaw, D. (2014). Network issue agendas on Twitter during the 2012 U.S. presidential election. Journal of Communication, 64(2), 296--316. \url{https://doi.org/10.1111/jcom.12089}
\end{enumerate}
}
\subsection*{Manuscripts Under Review}
{\small
\begin{enumerate}
\item Vargo, C. J., Dobolyi, D., \& Rahman, S. (2025, Oct 1). Inferring advertising demographics from behavioral data using a local LLM. Manuscript submitted for review to KDD (submitted Feb 1, 2026).
\item Vargo, C. J., \& Barrett, B. (2025, Sept 8). Who gets what ads? Tail-aware metrics and the composition of mobile political exposure. Letter of inquiry submitted to \textit{Journal of Quantitative Description: Digital Media} (submitted Jan 23, 2026).
\item Vargo, C. J., \& Hopp, T. (2025, Aug 7). Advertising, Not Legalization: Multi Level Evidence Linking Online Gambling Promotion to Problem Gambling in the United States. Manuscript submitted for review to Policy \& Internet (submitted Jan 6, 2026).
\item Vargo, C. J., \& Dobolyi, D. (2025, June 4). The hidden costs of misclassified content: Advancing press sustainability and advertising integrity through improved classification [Submitted to the Journal of Interactive Marketing; revise and resubmit; revision in process]
\end{enumerate}
}
\subsection*{Peer Reviewed, Edited Chapters}
{\small
\begin{enumerate}
\item Vargo, C., \& Hopp, T. (2018). Is Yik Yak a platform for political communication? Exploring college students' communication on an emergent social media platform. In N. J. Stroud \& S. McGregor (Eds.), Digital discussions: How big data informs political communication (pp. 144--165). New York, NY: Routledge.
\end{enumerate}
}
\subsection*{Edited Book Chapters}
{\small
\begin{enumerate}
\item Vargo, C., \& Guo, L. (2015). Exploring the network agenda setting model with big social data. In L. Guo \& M. McCombs (Eds.), The power of information networks: New directions for agenda setting (pp. 55--65). New York, NY: Routledge. \url{https://doi.org/10.4324/9781315726540}
\end{enumerate}
}
\subsection*{Refereed Conference Proceedings}
{\small
\begin{enumerate}
\item Vargo, C., Hopp, T., \& Gangadharbatla, H. (2025). From Ads to Addiction: The Role of Online Gambling Advertising Spend in Problem Gambling Search Trends over a Decade. Proceedings of the 58th Hawaii International Conference on System Sciences. Track: Digital and Social Media, Minitrack: Communication, Digital Conversation, and Media Technologies.
\item Vargo, C., Masullo, G., \& Hopp, T. (2024). Deciding to Delete Posts on Reddit: What Factors Influence Content Removal. Proceedings of the 57th Hawaii International Conference on System Sciences. Track: Digital and Social Media, Minitrack: Decision Making in Online Social Networks.
\item Vargo, C., Hopp, T., \& Agarwal, P. (2023). Inside a Social Media Brand Safety Algorithm: A Computational Investigation of Subreddits, Toxicity, and Advertising Inventory. Proceedings of the 2023 American Academy of Advertising Annual Conference. [With graduate advisee Pritha Agarwal]
\item Vargo, C., \& Hopp, T. (2023). Incivility on Popular Politics and News Subreddits: An Analysis of In-groups, Community Guidelines and Relationships with Social Media Engagement. Proceedings of the 56th Hawaii International Conference on System Sciences. Track: Digital and Social Media, Minitrack: Mediated Conversation. Maui, HI.
\item Gangadharbatla, H., Jones, V., Vargo, C., Ferrel, C., Owens, C. \& Li, H. (2019). The Role of Artificial Intelligence (AI) in Advertising. Proceedings of the 2019 American Academy of Advertising Annual Conference.
\item Vargo, C., \& Hopp, T. (2019). The effects of ad negativity on political digital advertising engagement: A computational case study of Russian Facebook and Instagram content. Proceedings of the 2019 American Academy of Advertising Annual Conference.
\item Hopp, T., \& Vargo, C. (2016). Does negative campaign advertising stimulate uncivil communication on social media? A big data analysis. Proceedings of the 2016 American Academy of Advertising Conference, 2016, 152--153.
\item Hester, J., \& Vargo, C. (2013). Social network sites and social media: A new research paradigm for strategic communication? Proceedings of the 2013 American Academy of Advertising, 2013, 177--178.
\end{enumerate}
}
\subsection*{Conference Papers}
{\small
\begin{enumerate}
\item Wang, Z., Guo, L., \& Vargo, C. (2025, August). Three worlds imagined through news: A cross-national analysis of country-based issue-ownership networks. Paper accepted for presentation in the International Communication Division at the Annual Conference of the Association for Education in Journalism and Mass Communication (AEJMC), San Francisco, CA, United States.
\item Barrett, B., \& Vargo, C. J. (2025, June 16). Divided by ads: Understanding partisan consumerism through online ad exposure. Paper presented in the ``High-Density: Strategies in Political Advertising and Campaigns'' session at the 75th Annual Conference of the International Communication Association, Denver, CO, United States.
\item Hopp, T., Ferrucci, P., Vargo, C., \& Mays, B. (2024). Mainstream News Media Trust, Countermedia Attendance, and Political Learning. Paper presented at AEJMC 2024 Annual Conference, Political Communication Division, Philadelphia, PA. [With graduate advisee B. Mays]
\item Vargo, C. \& Amazeen, M., (2020, August). The Competing ``Content Studio'' Agenda: A Large--scale Analysis of Sponsored Content in Elite U.S. Newspapers and Its Agenda Cutting Effect on Corporate News. Paper presented at AEJMC 2020 Annual Conference, Newspapers and Online News Division, Online.
\item Vargo, C., \& Liu, L. (2020, August). Audience ``agendamelding'' in the United States: assessing the relative influence of different types of online news on partisan agendas from 2015 to 2020. Paper presented at AEJMC 2020 Annual Conference, AEJMC 2020 Theory Colloquium on Agendamelding, Online.
\item Hopp, T., Ferrucci, P., Vargo, C., \& Liu, L. (2020, August). Is Facebook--Based Political Talk Associated with Political Knowledge? Paper presented at AEJMC 2020 Annual Conference, Political Communication Division, Online.
\item Vargo, C. J. (2020). Agendamelding at Scale: Assessing the Influence of Partisan, Emerging, Mainstream and Elite News Media on U.S. Audiences from 2015 to 2019. Panelist presentation at AEJMC 2020 Virtual Conference, session S066.
\item Amazeen, M., \& Vargo, C. (2019, August). Sharing native advertising on Twitter: Evidence of the inoculating influence of disclosures. Paper presented at AEJMC 2019 Annual Conference, Mass Communication and Society Division, Toronto, Canada.
\item Blankenship, J., \& Vargo, C. (2019, August). The effect of corporate media ownership on depth of local coverage and issue agendas: A computational case study of six Sinclair TV station websites. Paper presented at AEJMC 2019 Annual Conference, Electronic News Division, Toronto, Canada.
\item El--Toukhy, S., Vargo, C., Hopp, T., \& Choi, T. (2018, February). State--level demographics and tobacco--control correlates of smoking cessation behavioral change techniques on Twitter. Paper presented at the 2018 Society for Research on Nicotine and Tobacco Annual Meeting, Baltimore, MD.
\item Hopp, T., Ferrucci, P., \& Vargo, C. (2018, August). A citizen--based profile of fake news dissemination on Facebook. Paper presented at AEJMC 2018 Annual Conference, Political Communication Interest Group, Washington, DC.
\item Hopp, T., Ferrucci, P., \& Vargo, C. (2018, August). Social capital, civic engagement and identity: Exploring a model for political talk on Facebook. Paper presented at AEJMC 2018 Annual Conference, Political Communication Interest Group, Washington, DC.
\item Guo, L., \& Vargo, C. (2017, May). Who determines the global news agenda? A big data analysis of international news flow on the internet. Paper presented at ICA 2017, Mass Communication and Society Division (Top Faculty Paper for Journalism Studies Division), San Diego, CA.
\item Vargo, C., \& Guo, L. (2016, August). A network approach to intermedia agenda--setting: A big data analysis of traditional, partisan, and emerging online U.S. news. Paper presented at AEJMC 2016 Annual Conference, Online News Division, Minneapolis, MN.
\item Vargo, C., \& Hopp, T. (2016, September). Is Yik Yak a platform for political communication? Exploring college students' communication on an emergent social media platform. Digital discussions: How big data informs political communication. Paper presented at New Agendas Conference, Austin, TX.
\item Guo, L., \& Vargo, C. (2015, May). The power of ``issue ownership network'': A big--data analysis of the 2012 U.S. presidential election. Paper presented at ICA 2015, Political Communication Division, San Juan, Puerto Rico.
\item Hopp, T., \& Vargo, C. (2015, August). The effect of partisanship on changes in newspaper consumption: A longitudinal study (2008--2012). Paper presented at AEJMC 2015 Annual Conference, Mass Communication \& Society Division, San Francisco, CA.
\item Kim, Y., Gonzenbach, B., Vargo, C., \& Kim, Y. (2015, May). First and second levels of intermedia agenda setting: Political advertising, newspapers and Twitter postings during the 2012 U.S. presidential election. Paper presented at ICA 2015, CAT Division, San Juan, Puerto Rico.
\item Shaw, D., Vargo, C., Cole, R., \& Minooie, M. (2014, October). The emerging papyrus society: How we are using media to monitor civic life, find personal community and create private identity. Paper presented at the Arab--U.S. Association for Communication Educators Annual Conference, Irbid, Jordan.
\item Vargo, C., Guo, L., Shaw, D., \& McCombs, M. E. (2013, August). Network issue agendas on Twitter during the 2012 U.S. presidential election. Paper presented at AEJMC 2013 Annual Conference, Mass Communication Division, Washington, DC.
\item Vargo, C., \& Shaw, D. (2013, October). When motivated reasoning mediates agenda setting: Effects in the 2012 Republican presidential primaries. Paper presented at the Annual Conference of Media and the Public Sphere, Athens, Georgia.
\item Shaw, D., \& Vargo, C. (2012). The media and social stability: How audiences create personal community. Paper presented at the International Forum of Higher Education in Media and Communication, Beijing, China.
\item Vargo, C. (2012, March). Internet advertising \& interactive computer services: Liability \& immunity as provided by the Communications Decency Act. Paper presented at the 2012 Annual Southeast Colloquium of the AEJMC, Blacksburg, VA.
\item Vargo, C., \& Arguello, J. (2012, August). The unintended consequences of ``Moderate Mitt:'' The ideologies of Mitt Romney \& second‐level agenda setting. Paper presented at AEJMC 2012 Annual Conference, Chicago, IL.
\item Vargo, C. (2011, August). Twitter as public salience: An agenda--setting analysis. Paper presented at AEJMC 2011 Annual Conference, St. Louis, MO.
\end{enumerate}
}
\subsection*{Editorials (The Agenda Setting Journal)}
{\small
\begin{enumerate}
\item Vargo, C. J. (2021). Note from the editor: Online conference follow-up. The Agenda Setting Journal, 5(1), 1--7. \url{https://doi.org/10.1075/asj.00009.edi}
\item Vargo, C. J. (2020). Note from the editor. The Agenda Setting Journal, 4(2), 171--172. \url{https://doi.org/10.1075/asj.00007.edi}
\item Vargo, C. J. (2019). Note from the editor: Key advantages of submitting work to the agenda setting journal. The Agenda Setting Journal, 3(2), 103--105. \url{https://doi.org/10.1075/asj.00002.edi}
\item Vargo, C. J. (2019). Note from the editor. The Agenda Setting Journal, 3(1), 1--2. \url{https://doi.org/10.1075/asj.00001.edi}
\end{enumerate}
}
\subsection*{Invited Guest Lectures}
{\small
\begin{itemize}
\item Vargo, C. J. (2025, April 8). Brand safety and contextual suitability with machine learning applications. Invited guest lecture for Computational Advertising (graduate level), University of Illinois Urbana--Champaign.
\item Vargo, C. J. (2025, May 6). Doing new media research: Methods, strategy, and emerging directions. Invited guest lecture for ``New Media Research'' (graduate level), international Zoom presentation at Fudan University.
\item Vargo, C. J. (2024, September 24). Agenda setting theory: Foundations and modern applications. Invited guest lecture delivered via Zoom to incoming M.S. in Journalism students, Reed School of Media and Communications, West Virginia University.
\item Vargo, C. J. (2023, July 11). Computational social science across academia and industry. Invited guest lecture at the Summer Institute in Computational Social Science (SICSS--Beijing), Renmin University of China; international program featuring invited faculty from Northwestern University, Duke University, UC Davis, Purdue University, and others.
\item Vargo, C. J. (2020, October 28). Big data methods for communication research. Invited guest lecture delivered via Zoom to COMM 7020 (graduate level), School of Communication and Journalism, Auburn University.
\item Vargo, C. J. (2019, April 10). Invited research talk in the Information Science seminar series, Department of Information Science, University of Colorado Boulder.
\end{itemize}
}
\section{Awards}
\begin{itemize}
\item 2025, Association for Information Systems (AIS), International Conference on Information Systems (ICIS) 2025 Outstanding Reviewer.
\item 2023, Hawaii International Conference on System Sciences (HICSS-56), Best Paper Award (Minitrack: Mediated Conversation, Digital and Social Media Track) for paper, ``Incivility on Popular Politics and News Subreddits: An Analysis of In-groups, Community Guidelines and Relationships with Social Media Engagement.''
\item 2020, The University of Colorado Boulder, Provost's Faculty Achievement Award for Pre-Tenure Faculty. For the article: ``The agenda--setting power of fake news: A big data analysis of the online media landscape from 2014 to 2016.'' Awarded \$1,000.
\item 2018, JMCQ Outstanding Research Article of the Year Winner for: Vargo, C., \& Guo, L. (2017). Networks, big data, and intermedia agenda--setting: An analysis of traditional, partisan, and emerging online U.S. news. Journalism \& Mass Communication Quarterly, 94(4), 1031--1055.
\item 2017, Top Faculty Paper, Journalism Studies Division, ICA International Conference in San Diego, CA.
\end{itemize}
\section{Grant Applications \& Funds Raised}
{\small
\begin{itemize}
\item Vargo, C. (2023) Dissecting the Anatomy of 2024 Election Disinformation: A Graph-Based Approach. Submitted to neo4j upon request. \$164,143. Not funded.
\item Guo, L. \& Vargo, C. (2020) The agenda--setting power of fake news: The transfer of misinformation salience in a global, cross--platform information ecosystem. Submitted to Facebook. Grant duration: 2020 -- 2021. Grant amount: \$149,800. Submitted April 15th. Outcome: not funded.
\item Vargo, C., Hopp, T., \& Bradley, S. (2019). Socialcontext.ai -- An ai--powered social media content strategy tool built specifically for PR agencies. Submitted to the University of Colorado's Lab Venture Challenge, August 20th, 2019. Requested \$125,000. Finalist, but not funded.
\item Hopp, T., Ferrucci, P. \& Vargo, C. (2018) We the people: Fake news dissemination as a byproduct of citizens' civic practices. Arthur W. Page Center 'Fake News' Call for Grants, 2018. \$8,500. Funded in Full.
\item Hopp, T., Vargo, C., Ferrucci, P., Gondwe, G., Keegan, B., McDevitt, M., \& Skewes, E. (2018). A proposal to fund the Digital Citizenship Project. de Castro Research Reward at University of Colorado, College of Media Communication \& Information. \$24,934. Not Funded.
\item Vargo, C., \& Ferrucci, P. (2018) Social Listening Studio course proposal. Payden Teaching Grant at University of Colorado, College of Media Communication \& Information. \$7,500. Funded.
\item Vargo, C., \& Hopp, T. (2018). Twitter Health Grant: Discussion productivity potential. Twitter \$50,000. Submitted on April 12th. Selected as top 5 finalist amongst 239 submissions. Not Funded.
\item Vargo, C. (2016). National Institutes of Health (NIH) (award number/mechanism: \emph{[add grant/award ID here]}). Health disparity, tobacco use and socioeconomic factors: A longitudinal study of Twitter users across the United States. \$25,000. Funded.
\item MacCall, S. L., McMillan, D. J., Vargo, C. J., Bradley, S. B., \& Aversa, E. A. (2016). Knight Foundation's news challenge for libraries: How might libraries serve 21st century information needs?: Efficiency, integration, interoperability: A 21st century approach to organizing sports digital assets for all libraries -- pre--budget grant submission. Not funded.
\item Hopp, T., Vargo, C., \& Ames, B. (2015). Development of a scalable, multimodal identifier of malware propagation likelihood for NSF SaTC New Collaboration EAGER. Returned on March 27, 2015. Not selected to submit full proposal.
\item MacCall, S. L., Vargo, C. J., Bradley, S. B., \& Aversa, E. A. (2015). National Science Foundation -- Small Business Innovation Research (SBIR) Phase I Grant: Development of a novel digital asset organizing method in sports -- \$225,000 (\$74,925 sub--award to University of Alabama). Not funded.
\item Vargo, C. (2015). The University of Alabama -- Research Grants Committee Area B Award. Understanding how issues and attributes are created and evolve on social media using big data and Exponential Random Graph Models (ERGMs). \$6,000. Funded June 1, 2015.
\end{itemize}
}
\section{Teaching}
\subsection*{For Credit Coursera Specialization -- Text Analytics in Marketing}
{\small This three-credit sequence of courses was created for the University of Colorado's Master of Science in Data Science, a fully online graduate degree program offered via Coursera. Each course applies leading computer science methods to solve marketing problems using Python. Students receive conceptual overviews of methods, and dive into real-world datasets through instructor-led tutorials. Each course culminates with a major project. Last refreshed with new content Fall 2025.}\par
\vspace{0.5em}
\cvcourse{Supervised Text Classification for Marketing Analytics}{February 2022 -- Present}{This foundational course establishes comprehensive mastery of traditional text classification methodologies applied to authentic marketing datasets. Students develop advanced Python programming skills specifically tailored for text analytics, implement sophisticated keyword analysis and content analysis techniques with intercoder reliability validation, and master multiple machine learning algorithms including Support Vector Machines, Naive Bayes, Random Forest, and ensemble methods.}
\cvcourse{Advanced LLM Classification Systems and Production Deployment}{April 2022 -- Present}{This cutting-edge course provides comprehensive mastery of state-of-the-art Large Language Model classification systems for enterprise marketing applications. Students gain deep expertise in transformer architectures (GPT, BERT, LLaMA variants), advanced parameter-efficient fine-tuning techniques (LoRA, QLoRA, Unsloth), high-performance inference optimization (vLLM, quantization), and production-scale deployment strategies. The curriculum emphasizes hands-on implementation using industry-leading tools including Hugging Face Transformers, vLLM for inference acceleration, comprehensive evaluation frameworks, bias detection methodologies, and ethical AI governance for responsible deployment.}
\cvcourse{Topic Modeling and Network Analysis for Marketing Insights}{April 2022 -- Present}{This advanced course integrates cutting-edge unsupervised learning methodologies combining sophisticated topic modeling techniques with comprehensive network analysis for strategic marketing intelligence. Students master state-of-the-art topic discovery methods including BERTopic with transformer embeddings, Non-negative Matrix Factorization (NMF), Top2Vec, and advanced clustering algorithms applied to authentic Amazon product reviews and Twitter social media data.}
\subsection*{Coursera Specialization -- Digital Advertising Strategy}
{\small This three-course (non-credit) sequence is fully online, asynchronous and offered via Coursera. Students receive conceptual overviews of advertising platforms, and dive into assignments that exercise critical thinking about digital advertising strategy. Last refreshed Summer, 2024.}\par
\vspace{0.5em}
\cvcourse{Search Advertising}{December 2019 -- Present}{This course shows small businesses how to create and execute search campaigns on Google Ads Search (formally AdWords). Through an introductory overview, students are guided through the official Google Ads Search training materials where they will ultimately earn an official Google Ads Search Certification. Beyond Google materials, practical campaign creation and optimization best practices and exercises are provided by real experts.}
\cvcourse{Social Media Advertising}{December 2019 -- Present}{This course unpacks small business use cases of Facebook, Instagram and Twitter advertising. From basic campaigns, to advanced techniques including lookalike modeling and audience retargeting, this course shows how to effectively advertise on three major social media platforms.}
\cvcourse{Native Advertising}{2019 -- Present}{This course outlines a case study where a small travel startup used native advertising to drive hotel sales. Execution strategies for a successful, no--creative native campaigns are laid out, including: gathering existing news coverage, ethical content seeding, and content generation.}
\subsection*{Traditional Courses}
\cvcourse{Quantitative Methods (PhD level)}{Fall 2024, Fall 2025}{Introduces graduate students to quantitative social scientific research in the communication sciences. Topics covered will include -- but are not limited to -- the following: variables; research questions and hypotheses; causality; measurement; research design; descriptive statistics; and inferential statistics.}
\cvcourse{Social Media Listening}{Fall 2021, Fall 2024, Fall 2025}{This course covers the various applications of social listening as well as the ethics of using social media data. Students use cutting edge social media tools to tackle real-world social listening tasks. This class focuses on social media research activities, including analytics, listening and engagement. Social media metrics are also introduced and discussed.}
\cvcourse{Digital Advertising}{Fall 2018, Fall 2019, Fall 2020, Spring 2022, Spring 2023}{Part of the Leeds MS in Business Analytics, Marketing Track. Covers traditional and emerging digital advertising platforms. Students execute ads at an advanced level and use leading analytic tools to assess advertising performance. Core advertising platforms covered include search, display, social media, native advertising, mobile and programmatic.}
\cvcourse{Advanced Advertising Analytics}{Spring 2019, Spring 2020, Spring 2021}{This course tackles advanced advertising and marketing analytics through three advanced methods aimed at solving these problems: time series analysis, topic modeling via unsupervised machine learning, and network analysis. Includes a deep dive into the leading computer science methods aimed at solving these methods using Python and other open--source packages. Students walkthrough conceptual overviews of the methods, and dive in to real--world datasets through instructor--led tutorials.}
\cvcourse{Principles of Advertising}{Fall 2016, Spring 2017, Maymester 2017, Summer A 2017, Fall 2017}{Introduction to the foundations, nature, and practice of advertising. Commercial aspects of communications. Ethical, legal and social responsibility aspects of advertising.}
\cvcourse{Content Analytics}{Fall 2016, Spring 2017, Fall 2017, Fall 2018, Spring 2019, Fall 2019, Fall 2021, Spring 2022, Fall 2022, Spring 2023}{This class focuses on measuring and evaluating the effectiveness of strategic communication content (i.e., ads, webpages, and social media). It also provides an overview of metrics and best practices (e.g. KPIs) by medium.}
\cvcourse{Social Media Strategy}{Summer 2016}{Provides students with the practical, theoretical and analytical knowledge and skills required to successfully develop, monitor and execute digitally based and social media campaigns. Students will acquire a skill set based on the demands of current industry practice.}
\cvcourse{Independent Study}{Spring 2015, Fall 2015}{Brands \& Virality. Independent study with students interested in brands' use of social media and how that use leads to diffusion of messages.}
\cvcourse{Social Media Persuasion}{Fall 2015}{Master's level course on the practice of creating, writing, editing and producing persuasive communication for advertising and public relations.}
\cvcourse{Public Relations Writing}{Fall 2014, Spring 2015, Fall 2015, Spring 2016}{Theory and practice involved in creating messages, including planning, writing, editing, production and evaluation.}
\cvcourse{Design Application}{Fall 2010, Spring 2011, Summer 2014, Summer 2015}{An introductory skills course for advertising and public relations majors and minors using Adobe Photoshop, Illustrator and InDesign. Students develop skills to design logos, ads, brochures, newsletters and other basic APR documents.}
\cvcourse{Visual Communication}{Summer 2011, Spring 2015}{Theory, concepts and aesthetics applied to commercial persuasion in advertising and public relations. Particular attention is given to the application of course content in desktop publishing.}
\cvcourse{Advertising Media}{Spring 2014}{200--level course in media planning, buying and selling. Emphasized focus in buying and planning for social media.}
\cvcourse{Agenda--Setting Theory Seminar}{Spring 2013}{400--level course on the current and past works in agenda setting. The goal of the course is to have each student conduct a research study in agenda setting and complete it by semester's end. Co--instructed with Dr. Donald Shaw.}
\cvcourse{Introduction to Multimedia Design}{Summer 2012, Summer 2013}{Entry--level course in multimedia storytelling that includes modules on theory, the profession, design, content gathering, editing, programming, publishing and usability.}
\cvcourse{Creative Copy \& Communication}{Fall 2012}{Junior and senior--level course that reviews creative theory from social science research and applies thinking to advertising campaigns. Students brainstorm and execute multi--platform campaigns, while reviewing creative best practices.}
\section{Service}
\subsection*{Service to the University of Colorado Boulder}
\begin{itemize}
\item Faculty Representative, OIT Cabinet AI Steering Committee, Spring 2026 --- Present
\item Faculty Representative, OIT Cabinet OpenAI Steering Committee, Spring 2026 --- Present
\item Member, Boulder Faculty Affairs Campus Operations \& Resources Committee (CORC), Fall 2025 --- Present
\item Member, Strategic Facilities Visioning Committee, Fall 2018 --- 2024
\item Faculty Partner, Office of Academic and Learning Innovation, Fall 2018 --- Present
\end{itemize}
\subsection*{Service to the College of Media, Communication \& Information}
\begin{itemize}
\item Faculty Council Member, Fall 2021 --- Present
\item Personnel Committee Member, Fall 2021 --- Spring 2023
\item Director of Master of Science in Business Analytics (MSBA), Spring 2018 --- 2023
\item Co--organizer for Business Analytics Meetup Lecture Series, Fall 2018 --- 2022
\end{itemize}
\subsection*{Service to the Advertising, Public Relations and Media Design Department}
\begin{itemize}
\item Research roundtable presentation (30 minutes), CMCI APRD Journalism, University of Colorado Boulder, November 5, 2025. Demonstrated workflows for pulling, downloading, and extracting large-scale digital trace data using Meltwater, with examples for graduate students and junior researchers.
\item APRD/Journalism Research Colloquium Coordinator, AY 2025
\item Graduate Affairs Committee, AY 2023, AY 2024
\item Associate Chair of Graduate Studies, Fall 2022
\item Search Chair, 3 tenure-track lines in Strategic Advertising, Fall 2022
\item Member, Primary Unit Evaluation Committee, Fall 2021
\item PUEC mid-tenure review (Mia Wang), Spring 2026
\item Search Chair, Scholar in Residence of Digital Advertising Strategy, Fall 2021 --- Spring 2022
\item APRD Chair Search Committee Member, Fall 2019 --- Spring 2020
\item APRD Data Analytics Instructor Search Committee Chair, Spring 2019
\item Chair and Member of Executive Committee, Spring 2017 --- Fall 2017
\item Industry Partnerships Committee Member, Summer 2017 --- Spring 2018
\item AD + PR in Paris Summer 2018 Faculty--led Global Seminar Co--Director, Fall 2017
\item Advised Internships (Summer 2017: 7, Fall 2017: 2, Summer 2018: 3, Spring 2019: 2, Summer 2019: 1, Fall 2019: 1, Fall 2020: 1, Spring 2021: 1; Spring 2022: 1)
\item Directed Independent Studies (Fall, 2017: 1 Graduate, 1 Undergraduate; Spring, 2022: 1 Graduate)
\item PR Search Committee Member, 2016
\end{itemize}
\subsection*{Service to the Academy}
\begin{itemize}
\item Discussant, Communication Theory \& Methodology Scholar-to-Scholar (Poster) Refereed Paper Session, AEJMC 2020 Virtual Conference.
\item Guest Co--Editor, Special Issue: Agenda Manipulation in a Platformized Attention Economy: Building, Cutting, and the Governance of Salience, Frontiers in Political Science, Accepted Oct 2025. (with Lei Guo, Milad Minooie). [Call for papers announced, accepting submissions]
\item Guest Co--Editor, Special Issue: Revisiting Agenda-Setting Theory in the Age of AI, Communication and Change, Acknowledged in print. (with Lei Guo).
\item Editor, the Agenda Setting Journal, December 2018 --- 2024.
\item External Reviewer for Tenure \& Promotion, University of Southern California (USC) Annenberg School for Communication and Journalism, 2022 (1 case).
\item Guest Editor, the Agenda Setting Journal, Spring 2018 --- Summer 2018.
\end{itemize}
\subsection*{Dissertation Committees}
\begin{itemize}
\item Primary advisor (through November 7, 2024), Pritha Agarwal, University of Colorado Boulder, 2022 -- 2024
\item Primary advisor (Spring 2026 --- Present), Shahed Rahman, University of Colorado Boulder
\item Member, Yiyan Zhang, Boston University, Spring 2020 -- Summer 2021
\end{itemize}
\subsection*{Thesis Committees}
\begin{itemize}
\item Matt Brockman, University of Arizona, Spring 2020
\item Tom Arenberg, University of Alabama, Fall 2015 -- Fall 2017
\end{itemize}
\subsection*{Manuscript Reviewer}
{\small
\begin{multicols}{2}
\begin{itemize}
\item American Academy of Advertising (AAA), 2017, 2019, 2022
\item AEJMC National Conference, Ad Division, 2016, 2017, 2019, 2020
\item AEJMC National Conference, Newspaper and Online News Division, 2020, 2021
\item AEJMC National Conference, PR Division, 2015
\item Big Data \& Society, 2018
\item Chinese Journal of Communication, 2020
\item Communication Methods and Measures, 2019, 2020
\item Communication Research, 2020, 2021
\item Computers in Human Behavior, 2018
\item Emerald Studies in Media and Communication, 2016
\item Harvard Kennedy School (HKS) Misinformation Review, 2021
\item International Journal of Communication, 2016 -- 2022
\item International Journal of Press and Politics, 2015
\item International Journal of Public Opinion Research, 2021
\item Information, Communication \& Society, 2020
\item Journalism \& Mass Communication Quarterly, 2015 -- 2021
\item Journalism, 2015, 2021
\item Journalism Studies, 2019
\item Journal of Broadcasting \& Electronic Media, 2019, 2022
\item Journal of Communication, 2015 -- 2017, 2019 -- 2022
\item Journal of Information Technology \& Politics, 2015 -- 2016, 2018 -- 2023
\item Journal of Interactive Advertising, 2018 -- 2019
\item Journal of Marketing Communications, 2017
\item Journal of Public Relations Research, 2019
\item Mass Communication \& Society, 2017, 2019, 2020
\item New Media \& Society, 2015, 2017, 2019--2022
\item Political Communication, 2014, 2016, 2020 -- 2021
\item Social Media + Society, 2018
\item Social Science Computer Review, 2015, 2019 -- 2021
\end{itemize}
\end{multicols}
}
\section{Prior Professional Experience}
\begin{itemize}
\item Founder and CEO, socialcontext LLC (D.B.A. socialcontext.ai), June 2019 -- 2025
\item Digital and Social Media Creative, Porter Novelli, Atlanta, GA, May -- August 2010
\item Digital and Media Coordinator for the Digital Marketing Department, Epic Records/Sony Music, New York, NY, August -- December 2008
\item Postgraduate Internship, Digital Marketing Department, Columbia Records/Sony BMG, New York, NY, June -- August 2008
\item Public Relations Coordinator, ``On the Lot'' Television Show, FOX/Mark Burnett Productions/DreamWorks, Los Angeles, CA, May -- August 2007
\item Public Relations Intern, College of Engineering, The Pennsylvania State University, State College, PA, August 2007 -- May 2008
\end{itemize}

\end{document}